\begin{theorem}[The protocol under its composed system]
  \label{thm:aba-protocol-composed}
  \lean{PLTS.ABA.ABDY.ProtocolState}\lean{PLTS.ABA.ABDY.ProtocolRelation}\lean{PLTS.ABA.ABDY.protocolRelation_mk}\lean{PLTS.ABA.ABDY.protocolRelation_init}\lean{PLTS.ABA.ABDY.protocolSimulation}\lean{PLTS.ABA.ABDY.protocol_composed}\leanok
  \uses{def:aba-composed, def:aba-programs, def:aba-network, def:aba-protocol, def:forwardsim, def:coupling, def:achievable, thm:forward-sound}
  For all parameters $P$, the protocol is forward-simulated by its composed system along the Dirac lift of the relation $\mathrm{ProtocolRelation}(P)$, and hence
  \[
    \operatorname{ATD}(\abdy{\mathrm{protocol}}(P)) \;\subseteq\;
    \operatorname{ATD}(\abdy{\mathrm{composed}}(P)).
  \]
  The two presentations hold the same round data indexed two ways, by process at the protocol and by round at the composed system.
  $\mathrm{ProtocolRelation}(P)$ has eight conjuncts.
  The first six relate the coordinates: the round loops agree, replacement flags included (D23); the common coin is the same component in both systems; the ABA network is the network's DECIDED sets beside its corrupted set and its started rounds; the rounds in the map of the composed family are the started rounds (D37); the network state of round $r$ is the network's round-$r$ sent sets beside the same corrupted set and its ghost for that round, the round's bound bit; and the entry of process $j$ in the instance of round $r$ is the variables of round $r$ that $j$ holds, a round $j$ has not touched reading as the initial round variables in both systems.
  The last two are invariants of the protocol state: at a round outside the started rounds no process has sent, the ghost holds no bit and no process holds variables; and every process of the corrupted set has a replaced program.
  A composed state related to a state of the protocol is therefore determined by it, up to the corruptions its family has recorded, which the protocol state does not hold.
  The determination is not injective: no composed state carries a termination flag, so two states of the protocol differing only in which processes have terminated are related to the same composed state (D22).
  Headline: \leandecl{PLTS.ABA.ABDY.protocol_composed}.
\end{theorem}

\begin{proof}\leanok
  \uses{def:aba-composed, def:aba-programs, def:aba-network, def:aba-protocol, def:forwardsim, def:coupling, def:achievable, thm:forward-sound}
  The initial states are related, \leandecl{PLTS.ABA.ABDY.protocolRelation_init}: everything is initial, no round is started, each process's map has no entry at any round and $\mathrm{roundVariables}\,r$ is the initial round variables there, which is the entry the composed initial state carries.
  For the step condition, take a state of the protocol, a composed state related to it, and a transition out of that protocol state.
  Every transition of the protocol is a product of Dirac factors beside the common coin's successor distribution, and so is the composed transition that matches it.
  The two presentations carry the same coin, so the coupling is functional in the coin coordinate and the matched outcomes differ in no coordinate the relation constrains.
  The transitions are matched by label class, each by the composed transition carrying the same label, by a silent transition of one round instance, or by the empty run.
  A graded-agreement call at an unreplaced program records the called round among the started rounds, and the family enters that round in its map on the same call.
  Every other transition at round $r$ is guarded by the started rounds or reads variables of round $r$ that hold an input, a message sent at round $r$ or the ghost of round $r$, so $r$ is a started round and a round of the family's map, where the family takes that round's silent transitions.
  A round's transition writes the variables of the round its label tags and no other, and the round instance of that round writes the entry of the acting process and no other, so the entry conjunct is preserved entry by entry.
  The round's first return writes the network's ghost and the instance's bound bit alike, every other round's ghost unchanged, and the network's announcement of the bound bit is matched by the instance's announcement, whose guard reads the same bit.
  A corrupted process's graded-agreement call leaves the protocol's state where it stands and is matched by the empty run; so are a round delivery at a replaced or terminated program and the termination of a process, which writes a flag the relation does not read.
  A corrupted process's coin call is matched by the composed system's own, which the ABA network outputs under the same two guards (D11, D37), and the replaced program's return by the replaced round loop's self-loop beside the ABA network's unguarded transition (D23).
  Two classes of composed transition have no protocol counterpart, which is what makes the inclusion one-directional: a round instance records a corrupted process's call, and the round transitions of a round instance carry no termination guard, so the instance takes a send or a delivery at a process the protocol has terminated.
  Soundness (Theorem~\ref{thm:forward-sound}) turns the simulation \leandecl{PLTS.ABA.ABDY.protocolSimulation} into the stated inclusion.
\end{proof}
